\section{The Fork: Recoil or Consent}

The metaphysical judgment and the spiritual reception of it must be distinguished. A contemplative experience cannot repair an invalid inference, and no particular sensation is promised to one who understands the argument. Yet patristic and mystical writers describe the interval in which a truth already known becomes difficult to hold at a distance.

\subsection{Arrest before response}

Augustine gives that interval an exact form. In \emph{On the Trinity} VIII.2.3, he asks the reader to attend to Truth itself. The mind receives something like a first flash, cannot sustain the gaze, and slips back into familiar earthly images. His subject there is not a proof of present conservation. His phenomenology is nevertheless apt: intellectual contact precedes the retreat by which the mind recovers its accustomed standpoint.

The \emph{Confessions} makes the response more personal. Before the immutable Light Augustine knows himself as made (VII.10.16), and later knows his changeable mind as judging mutable things by unchangeable truth (VII.17.23); before eternal Wisdom he both shudders and burns, \latin{inhorresco et inardesco} (XI.9.11). Trembling is not yet refusal. Disproportion may be felt at the same moment that likeness awakens desire. \emph{Recoil}, in the stricter sense used here, is the movement that interprets the disproportion as a reason to withdraw.

\begin{studybox}{A terse taxonomy}
\begin{description}[leftmargin=9.8em,style=nextline,itemsep=0.1em,topsep=0.15em]
\item[Apprehension] The intellect judges asymmetric dependence to be true.
\item[Arrest] The judgment turns upon the knower: \emph{this is true of me}.
\item[Recoil] The self withdraws as though dependence threatened its reality.
\item[Abasement] The self relinquishes false elevation and takes its creaturely place.
\item[Consent] The will freely agrees with the truth of creaturehood.
\item[Humility] That truth becomes a stable disposition.
\item[Magnanimity] Received goods are accepted as real powers to be exercised.
\item[Self-forgetfulness] Attention ceases to make the self its own measure.
\item[Adoration] The whole person renders God the reverence due to the first principle.
\end{description}
\end{studybox}

These are distinctions, not a guaranteed sequence of feelings or a ladder of mystical states. Dependence is true before consent; consent changes the person's relation to the truth, not the truth's ontological force.

\subsection{Proper abasement}

Aquinas prevents humility from becoming self-contempt by pairing it with magnanimity. A difficult good both attracts and occasions retreat. Humility restrains the appetite from reaching beyond right reason; magnanimity strengthens it against despair and urges it toward genuinely great goods according to right reason. The virtues are compatible because both preserve truth (ST II--II, q.~161, a.~1, especially ad 3; q.~129, a.~3, ad 4).

This pairing is decisive. Pride converts received good into an occasion for inordinate self-exaltation, sometimes treating it as though it came from oneself. False abasement denies that a good was given. Humility confesses, \emph{I am not from myself}; magnanimity adds, \emph{what I have received is truly mine to exercise}. Neither sentence survives without the other.

Aquinas's account of adoration makes the bodily language of abasement exact. Interior adoration is devotion and reverence; kneeling or prostration signifies that inward act. In prostration, he says, ``we profess that we are nothing of ourselves'' (ST II--II, q.~84, a.~2, ad 2). The claim is not \emph{I am nothing}. It is \emph{I am, but not from myself}. Devotion, the ready will to give oneself to divine service, supplies the volitional center of consent (q.~82, a.~1).

The \emph{Catechism} draws the same consequence. Adoration acknowledges God as Creator and ``infinite and merciful Love''; it confesses the creature who would not exist but for God, yet sets the worshiper free from turning in upon himself (CCC 2096--2097). C.~S. Lewis makes the attentional point in modern terms: genuine humility ends habitual self-occupation (\emph{Mere Christianity}, III.8). Proper abasement therefore begins with self-knowledge and ends by ceasing to make the self the horizon.

\subsection{The mystics: little, real, and loved}

Julian of Norwich holds littleness and goodness together. In chapters 5--6 of the long text of her \emph{Revelations of Divine Love}, the created whole appears as a tiny thing, about the size of a hazelnut. It lasts because God loves it: made, loved, and kept. Beholding the Maker makes the soul seem less in its own sight, but the result is reverent dread, true meekness, and charity rather than disgust with creaturehood. The image does no independent work as a proof. Its force is to make received being apprehensible at once.

Yet universal dependence is not sanctifying union. In \emph{Ascent of Mount Carmel} II.5.3--7, John of the Cross distinguishes God's substantial presence, by which every creature is preserved in being, from transformative union by likeness of love. A sinner does not cease to depend ontologically on God; a saint does not cease to be a creature when united to him.

Pseudo-Dionysius then names love's movement beyond self-enclosure. In \emph{Divine Names} IV.13--14, love is ecstatic: it carries the lover toward the beloved and returns toward the Good from which it proceeds. This is not annihilation of the created person. It is the redirection of attention and will away from self-possession as a final end.

Bernard of Clairvaux supplies the free answer. In \emph{Sermons on the Song of Songs} 83.3--6, the soul is conformed to the Word through the will, loving because it has been loved. Its return is whole but never equal---\latin{etsi non ex aequo}. The lover is not Love itself; the thirsty one is not the fountain. Yet love may be wholly the creature's act while continually received from its source.
